\section{Evaluation}
\label{sec:evaluation}

We answer four research questions.
\begin{description}
  \item[\rqref{1} Effectiveness.] How many previously unknown bugs does \tool find in DuckDB~2.0.0, of which kinds, and
        how do maintainers respond?
  \item[\rqref{2} Methods.] What does each of the four methods contribute, and at what cost?
  \item[\rqref{3} Judge.] Does the engine judge separate generators that exercise their target from generators that
        only claim to?
  \item[\rqref{4} Loop.] How efficient is the unattended rule-targeted loop, from LLM calls to filed issues?
\end{description}

\paragraph{Setup.} All runs of Methods A, B and D and the Claude arm of Method C used one laptop: an Intel Core
i3-1315U (8 hardware threads), 7.4\,GB of RAM under WSL2 (Linux 6.18), Python~3.14.4 and the released DuckDB~2.0.0
(source tag \texttt{v2.0.0}); Method A on DuckDB~1.5.6 and~2.0.0 used our own sanitizer builds. The LLM of Method D was
Claude Code~2.1.323 with Claude Opus for the first round of each target and Claude Sonnet for repair rounds, on a
personal subscription. Round~2 of Method C ran on a team member's machine with Gemini~2.5~Pro. Workers ran one at a
time with a DuckDB memory limit, because larger settings froze the machine. Every number in this section is produced
by \texttt{paper/scripts/collect\_data.py} and \texttt{collect\_extra.py} from run logs and from the GitHub state of our
issues on 13~November~2026; the outputs are in \texttt{paper/data}.

\subsection{RQ1: Bugs Found}
\label{sec:eval-bugs}

%% generated by paper/scripts/collect_data.py from data/issues.csv
\begin{table*}[t]
\caption{All 16 issues we filed on DuckDB. \emph{Method}: the method of Section~\ref{sec:approach} that found it. \emph{Kind}: wrong result (silent), internal error (an assertion or INTERNAL Error, in two cases also invalidating the database), or spurious error (a valid query fails). \emph{Status} on 2026-11-13; fix PRs are by maintainers or other contributors, except \#28547 (ours).}
\label{tab:issues}
\footnotesize\centering
\begin{tabular}{llllp{4.1cm}}
\toprule
Issue & Method & Kind & What goes wrong & Status \\
\midrule
\#27759 & B & internal error & FOR EACH ROW trigger into a table with a BEFORE INSERT trigger & fixed (PR \#28698) \\
\#27850 & C & internal error & identical EXISTS blocks under UNION ALL (delim\_join\_as\_cte=false) & fixed (PR \#28144) \\
\#27904 & C & internal error & CROSS JOIN LATERAL with a nested LATERAL count(*) & open, fix PR \#27975 \\
\#28052 & C & wrong result & negated DOUBLE compared with a NaN constant & fixed (PR \#28071) \\
\#28527 & D (session) & wrong result & date\_trunc with a sub-day part on a DATE column & open, fix PR \#28528 \\
\#28531 & D (session) & spurious error & CAST(timestamp AS DATE) = date at the TIMESTAMP range limits & open, fix PR \#28547 \\
\#28532 & D (session) & internal error & TIMESTAMP\_MS to TIMESTAMP\_S cast at range end invalidates the database & open, fix PR \#28539 \\
\#28534 & D (session) & spurious error & date\_trunc comparison near the end of the TIMESTAMP range & open \\
\#28535 & D (session) & internal error & UBIGINT constant above BIGINT max in x + c1 = c2 & open, fix PR \#28540 \\
\#28536 & D (loop) & wrong result & regex empty group (?:) at the pattern edge treated as an anchor & open, fix PR \#28538 \\
\#28541 & D (loop) & wrong result & UHUGEINT constant above HUGEINT max wraps in x + c1 = c2 & open, fix PR \#28549 \\
\#28542 & D (loop) & wrong result & TRY\_CAST(date AS TIMESTAMP) comparison returns false instead of NULL & open \\
\#28559 & D (loop) & wrong result & top-N window elimination returns the other join side for 0.0 / -0.0 & open, fix PR \#28666, \#28728 \\
\#28562 & D (loop) & wrong result & contains(list, '') returns true for lists without an empty string & fixed (PR \#28563) \\
\#28565 & D (loop) & spurious error & CASE guard around a STRUCT field cast ignored & open, reproduced \\
\#28566 & D (loop) & wrong result & both sides of a DOUBLE[] join key merged for [0.0] and [-0.0] & open, fix PR \#28728 \\
\bottomrule
\end{tabular}
\end{table*}


Table~\ref{tab:issues} lists the 16 issues we filed between 23~October and 9~November~2026; a team member filed one
(\#27904) and another team member filed the rest. Each issue has a minimized reproducer, the affected versions, and the
result with the responsible switch turned off. By kind, 8 are silent wrong results, 5 are internal errors, and 3 are
spurious errors on valid queries. Two of the internal errors are severe: \#28532 invalidates the whole database, so
every later query in the session fails, and \#27759 aborts a trigger cascade.

Maintainers have fixed 4 issues: \#28052 (PR \#28071), \#27850 (PR \#28144), \#28562 (PR \#28563) and \#27759
(PR \#28698, which also fixes a related wrong result; DuckDB's lead developer closed the issue). For 9
more issues a fix pull request is open; 8 of them were written by maintainers or other contributors after reading our
report, and one (\#28547 for \#28531) is ours. DuckDB's bot labelled 11 of the 16 issues as \emph{reproduced}. None was
closed as a duplicate or as expected behaviour.

The bugs cluster in one family: constants at the limits of their type inside optimizer rewrites. The rewrite builds a
new value (a bound, a moved constant, a cast) and that value overflows, wraps or loses a NaN or a sign. In \#28541, a UHUGEINT constant above the HUGEINT maximum silently wraps when it is moved across \texttt{=}. In \#28534, a \texttt{date\_trunc} bound lies one step past the last valid timestamp. In \#28559 and \#28566, the signed zero makes \texttt{[0.0] = [-0.0]} true although the two values print differently.

\begin{observation}
\tool found 16 previously unknown bugs in DuckDB~2.0.0, half of them silent wrong results. Maintainers fixed 4 within three weeks and fix pull requests exist for 9 more; none was rejected.
\end{observation}

\subsection{RQ2: What Each Method Contributes}
\label{sec:eval-methods}

\begin{table}[t]
\caption{What each method ran and what it found. \emph{Work}: tests executed or LLM calls spent. \emph{New}: issues
filed that were previously unknown. Method A ran first on DuckDB~1.5.6 and later on~2.0.0; all its real findings were
already known upstream.}
\label{tab:methods}
\footnotesize\centering
\begin{tabular}{lp{4.4cm}p{4.6cm}rr}
\toprule
Method & Input and LLM & Work & New & Fixed \\
\midrule
A classic & none & 70{,}980 SQLsmith queries (92.3\% valid); 11{,}235 differential queries; 14{,}000 sanitizer runs & 0 & 0 \\
B features & hand-written fuzzers for new 2.0 features & trigger, window, join, bottom-up and delim-join fuzzers & 1 & 1 \\
C hints + LLM & static and new-code hints, feature cards, dossiers; 1.5B local model, Gemini~2.5~Pro, Claude &
  Gemini Round~2: 12 generators, 36{,}600 cases, 122{,}590 valid queries; Claude arm & 3 & 2 \\
D rule loop & rule cards; Claude Opus and Sonnet & interactive session; loop runs \texttt{smoke1} (2 calls) and
  \texttt{hunt1} (40 calls) & 12 & 1 \\
\midrule
Total & & & 16 & 4 \\
\bottomrule
\end{tabular}
\end{table}

Table~\ref{tab:methods} compares the methods. \textbf{Method A} found real failures but no new bug. A differential
sweep of 11{,}235 queries over DuckDB's own tests on DuckDB~1.5.6 surfaced two crashes, both already reported and fixed
upstream; a sanitizer hunt over 14{,}000 runs produced 109 hits with 22 distinct signatures, all of them known bugs or
expected behaviour; and a 45-minute SQLsmith campaign of 70{,}980 queries on the 2.0.0 sanitizer build produced 46 raw
anomalies and no new bug. This is the gap that motivated the other methods: generic fuzzing re-finds what upstream's
own fuzzers already find.

\textbf{Method B} found our first new bug: the trigger fuzzer found \#27759, an internal error when a
\texttt{FOR EACH ROW} trigger inserts into a table that has a \texttt{BEFORE INSERT} trigger. Triggers were new in
DuckDB~2.0 and rarely appeared in upstream reports.

\textbf{Method C} found three bugs (\#27850, \#27904, \#28052), two of them now fixed. Its results depend strongly on
the model. In the pilot with a local 1.5B-parameter model, 72 of 126 guided queries were valid (57\%) and none of 10
targeted code regions was reached. Gemini~2.5~Pro in Round~2 wrote 12 generators that passed the audit and ran
122{,}590 valid queries without a new bug; its four candidates were two sanitizer-only failures already fixed upstream
and two that did not reproduce. The same round shows the cost of directing: the directed generators covered 17.89\% of
DuckDB's 160{,}389 instrumented lines, against 38.64\% for the baseline corpus, because each generator stays close to
one target.

\textbf{Method D} found 12 of the 16 bugs: 5 in an interactive session in which Claude read the optimizer source and its fix history with a person in the loop, and 7 in the unattended loop. Rule cards automate what this session did by hand. These 7 came from 42 LLM calls (Section~\ref{sec:eval-loop}).

\begin{observation}
The methods are complementary. Generic fuzzing found no new bug, a feature fuzzer found the first one, and directing an
LLM at single rules found 12 of 16. A larger model and a narrower target mattered more than more tests: Method C ran
more than 120{,}000 queries in one round without a new bug, and Method D found 6 bugs with 40 calls.
\end{observation}

\subsection{RQ3: Does the Engine Judge Work?}
\label{sec:eval-judge}

\begin{table}[t]
\caption{Controls of the judge and the triage. A control must reach its expected verdict after every change to them.
\emph{Rule}/\emph{Expr.}: rule-level and expression-level reach rate. \emph{Re-found}: the issue that triage classified the
control's hits as a rediscovery of.}
\label{tab:controls}
\footnotesize\centering
\begin{tabular}{llrrrl}
\toprule
Control & Judge verdict & Rule & Expr. & Hits & Re-found \\
\midrule
cheat (names rule, never fires it) & rejected & 0.00 & 0.23 & 0 & -- \\
positive (\texttt{date\_trunc}) & accepted & 0.33 & 0.55 & 8 & \#28527 \\
session generator, arithmetic & rejected (diversity) & 0.59 & 0.59 & 8 & \#28535 \\
session generator, cast comparison & rejected (diversity) & 0.11 & 0.40 & 12 & \#28532 \\
session generator, \texttt{date\_trunc} edges & rejected (diversity) & 0.22 & 0.56 & 6 & \#28534 \\
session generator, timestamp cast & rejected (diversity) & 0.35 & 0.58 & 2 & \#28531 \\
\bottomrule
\end{tabular}
\end{table}

The judge exists because an LLM reports progress it did not make. In the first unattended Method~C run, 10 of 11
generators written by an LLM agent produced exactly 8 queries per case and 5 to 12 distinct plans, and the run's only
reported bug was our own \#27759 again. Our first judge counted SQL keywords instead of operators, which a red-team
review showed could be satisfied without testing anything; we rebuilt it on operator trees from
\texttt{EXPLAIN (FORMAT JSON)}.

Table~\ref{tab:controls} shows the controls we run after every change. The negative control names the target rule but
never gives it a bare column, so constant folding removes the call before the rule sees it. Its expression-level reach
is 0.23, but its rule-level reach is 0.00, and the judge rejects it. The positive control is accepted and its hits are
classified as rediscoveries of \#28527. The four session generators that found \#28531, \#28532, \#28534 and \#28535
still re-find exactly their own bug, so the triage and the registry work end to end.

The same table also shows a limit of the judge. All four session generators fail the diversity floor: they are narrow
and deep, with few query skeletons, yet each found a real bug. For this reason Method D triages the hits of every
generator. In \texttt{hunt1}, the judge accepted 23 of 32 rounds, and all 9 candidates came from accepted rounds; the rejected rounds produced 128 raw hits and no candidate.

\begin{observation}
Measured rule-level reach rejects a generator that only names its target, and controls that re-find our own bugs show
the pipeline still works after every change. The diversity floor is a measure of coverage, not of bug-finding power:
narrow generators can find bugs, so their hits must still be triaged.
\end{observation}

\subsection{RQ4: The Unattended Loop}
\label{sec:eval-loop}

%% generated by paper/scripts/collect_data.py from data/hunt_units.csv and the hunt1 ledger
\begin{table}[t]
\caption{The automated loop (Method D, run \texttt{hunt1}): most targets give nothing, a few give bugs. \emph{Units}: generator rounds run. \emph{OK}: rounds the engine judge accepted. \emph{Reach}: best fraction of valid queries whose plan changes when the target is switched off. \emph{Hits}: raw differences before triage. \emph{Cand.}: hits that survived triage. \emph{Issues}: filed upstream.}
\label{tab:hunt1}
\footnotesize\centering
\begin{tabular}{lrrrrrl}
\toprule
Target & Units & OK & Reach & Hits & Cand. & Issues \\
\midrule
\texttt{remove\_unused\_columns} & 4 & 4 & 1.00 & 95 & 2 & 28565, 28566 \\
\texttt{date\_trunc\_simplification} & 2 & 1 & 0.61 & 0 & 0 & -- \\
\texttt{move\_constants} & 2 & 2 & 0.98 & 98 & 1 & 28541 \\
\texttt{comparison\_simplification} & 2 & 1 & 0.31 & 189 & 1 & 28542 \\
\texttt{monotone\_preimage} & 3 & 3 & 0.80 & 10 & 0 & -- \\
\texttt{timestamp\_comparison} & 3 & 3 & 0.77 & 3 & 0 & -- \\
\texttt{string\_prefix} & 3 & 2 & 0.64 & 1 & 0 & -- \\
\texttt{topn\_window\_elimination} & 3 & 3 & 0.97 & 108 & 4 & 28559 \\
\texttt{in\_clause\_simplification\_rule} & 3 & 1 & 0.37 & 60 & 0 & -- \\
\texttt{compressed\_materialization} & 1 & 0 & 0.59 & 0 & 0 & -- \\
\texttt{empty\_needle\_removal} & 3 & 2 & 0.63 & 31 & 1 & 28562 \\
\texttt{list\_comprehension\_rewrite} & 3 & 1 & 0.60 & 2 & 0 & -- \\
\midrule
Total & 32 & 23 & & 597 & 9 & 6 \\
\bottomrule
\end{tabular}
\end{table}


Table~\ref{tab:hunt1} shows the run \texttt{hunt1}. It spent 40 LLM calls on 13 targets (12 of them ran at least one
round), ran 32 generator rounds with a median valid rate of 95\%, and collected 597 raw hits. Triage looked at 138 hits, at most ten per round. It classified 56 as duplicates within the run and 54 as already fixed on the development build. Another 10 were low priority (the optimizer only hid an error), 8 were rediscoveries of known issues and 1 was an artifact. The remaining 9 were new candidates. A person read the 9 candidates: four were the same top-N bug, one was covered by an open upstream
pull request, and six became issues. A short earlier run (\texttt{smoke1}, 2 calls) gave one more (\#28536). The run was started seven times; after a machine sleep, failed LLM calls and a usage limit, each start resumed from the ledger without repeating a completed round.

Yield is uneven. Five of the 12 targets produced all the candidates, and they share a trait: the rule builds a new
value, or it replaces one column with another that compares as equal. Rules whose queries rarely reach them
(\texttt{in\_clause\_simplification\_rule}, \texttt{compressed\_materialization}) produced nothing. The triage step saved much work by dropping 54 hits whose bug was already fixed on the development build, for example top-N hits fixed by PR \#28405 two days before the run.

\begin{observation}
The loop turned 40 LLM calls into 6 new issues, about one issue per 7 calls, with a person reading only 9
candidates out of 597 raw hits. Yield comes from a few targets that build values, so rotating targets helps more than
longer runs on one target.
\end{observation}
